你这个文档，最好不要写成“研究想法记录”，而要写成一份 **开发推进规范**。

它的作用只有一个：
**把“现在该做什么、做到什么算过关、不过关怎么办、Git 怎么推进”全部写死。**

按照你现在这份材料，最容易在实现里出问题的正是这些地方：离散残差怎么定、震源附近怎么 mask、PML 内介质怎么处理、哪些量能直接差分哪些不能直接差分。也就是说，文档的核心不是再讲理论，而是把这些实现决策冻结下来。 

我建议你直接在仓库里放一个文件：

`docs/DEV_PLAN.md`

然后按下面这个结构写。

---

# 你这份文档应该怎么写

## 1. 文档首页先写清楚三件事

### 1.1 项目目标

一句话说明你现在这个仓库要做什么。

例如：

> 本项目旨在实现一个面向高波数 Helmholtz 方程的分阶段原型系统，先完成 2D 单源、双精度、预处理与物理残差验证，再逐步扩展到神经算子训练与多波数泛化。

### 1.2 当前阶段目标

不要写太大。第一阶段只写能落地的。

例如：

> 当前阶段只做 2D、单源、标量声学、规则网格、双精度、单机 PyTorch 原型。
> 暂不处理 3D、多源批处理、分布式训练、大规模 benchmark。

### 1.3 开发原则

这一段非常重要，建议固定写：

* 一次只推进一个里程碑
* 上一步未通过检查，不进入下一步
* 所有公式先冻结，再编码
* 每个里程碑必须有最小测试
* 每次合并到主分支前，必须更新文档和实验记录
* 遇到理论与实现冲突，优先修改文档，不直接改代码绕过去

---

## 2. 写“范围冻结”

你这个项目最怕一开始范围太大。
所以文档里必须单开一节叫：

## Scope Freeze / 范围冻结

建议直接写成这样：

### 本版本包含

* 2D Helmholtz
* 单个点源
* 规则笛卡尔网格
* 复数场用实部/虚部双通道表示
* 因子化 Eikonal 预处理
* 背景场 `u0`
* 走时 `tau`
* `grad_tau`、`lap_tau` 的安全重组
* PML 复拉伸
* 固定卷积核差分
* 前向网络骨架
* PDE residual 检查

### 本版本不包含

* 3D
* 多源联合训练
* 自动微分求二阶空间导数
* 端到端大规模训练
* 论文级 benchmark
* 工程部署优化

这一节的意义是：以后你如果想临时加内容，先改文档，不要直接改代码。

---

## 3. 写“关键实现决策”

这部分是整份文档最值钱的地方。
不要怕短，但一定要硬。

建议列成“冻结决策”：

## Frozen Decisions / 冻结决策

### D1. 复数表示

* 采用实部/虚部双通道，不直接依赖复杂 dtype 做核心网络输入输出
* 所有物理张量接口统一为 `[B, C, H, W]`

### D2. 空间导数实现

* 不使用 `autograd` 计算二阶空间导数
* 一阶、二阶导数统一使用固定卷积核

### D3. 走时计算

* 不直接离散展开后的因子化方程
* 走时求解采用因果一致的 Godunov/FSM 路线
* `tau = tau0 * alpha`

### D4. `lap_tau` 的构造

* 不直接对 `tau` 做二阶差分
* 使用链式法则重组：
  `Δτ = αΔτ0 + 2∇τ0·∇α + τ0Δα`

### D5. 震源点处理

* PDE residual 在震源邻域采用 punctured-domain mask
* mask 半径在当前版本固定，例如 `1.5h`

### D6. PML 内介质处理

* 当前版本中，PML 层内介质退化为背景介质
* 因而等效体源在 PML 内视为 0

### D7. 当前训练策略

* 当前阶段不启动正式训练
* 先完成前向、残差、shape、数值稳定性检查

这些决策其实正好对应你材料里已经暴露出的几个高风险点。 

---

## 4. 写“仓库结构”

这一节可以非常实用。比如：

```text
repo/
├─ docs/
│  ├─ DEV_PLAN.md
│  ├─ THEORY_TO_CODE.md
│  └─ CHANGELOG.md
├─ configs/
│  ├─ base.yaml
│  ├─ debug.yaml
│  └─ m1_grid.yaml
├─ src/
│  ├─ core/
│  │  ├─ grid.py
│  │  ├─ medium.py
│  │  ├─ source.py
│  │  └─ complex_utils.py
│  ├─ physics/
│  │  ├─ background_field.py
│  │  ├─ factored_eikonal.py
│  │  ├─ tau_reconstruct.py
│  │  ├─ pml.py
│  │  ├─ diff_ops.py
│  │  └─ residual.py
│  ├─ models/
│  │  ├─ spectral_block.py
│  │  └─ nsno.py
│  ├─ train/
│  │  ├─ losses.py
│  │  └─ trainer.py
│  └─ tests/
│     ├─ test_grid.py
│     ├─ test_pml.py
│     ├─ test_background.py
│     ├─ test_tau.py
│     └─ test_residual.py
├─ notebooks/
├─ scripts/
└─ README.md
```

---

## 5. 核心部分：写“里程碑推进表”

这部分决定你能不能一步一步稳着写。

你不要写“先写预处理，再写网络”这种太空的话。
你要写成下面这种格式：

---

## Milestones / 里程碑

### M0. 建仓与最小骨架

**目标**
完成仓库初始化、目录、配置系统、基础 README。

**输出**

* 可运行的 Python 环境
* 目录结构固定
* `base.yaml`
* 空仓测试通过

**检查项**

* 能安装依赖
* `pytest` 能跑
* 配置能读
* 日志能输出

**通过标准**

* 新机器克隆后 10 分钟内可启动
* 无手工改路径

**Git**

* 分支：`feat/m0-bootstrap`
* 合并条件：目录、环境、README、基础测试齐全

---

### M1. 网格、介质、源项

**目标**
实现规则网格、介质场、点源坐标和基础数据结构。

**输出**

* `Grid2D`
* `Medium2D`
* `PointSource`
* 距离场 `r(x)`

**检查项**

* 网格 shape 正确
* 坐标系定义统一
* 源点索引和物理坐标一致
* 边界索引无错位

**通过标准**

* 3 个手写 case 都能对上
* 可视化正常

**Git**

* 分支：`feat/m1-grid-source`

---

### M2. 背景场与解析先验

**目标**
实现 `s0`、`tau0`、`u0`。

**输出**

* `s0 = s(xs)`
* `tau0`
* 2D/当前版本背景场 `u0`
* 震源点安全值处理

**检查项**

* 远离震源时趋势合理
* 震源点无 NaN/Inf 泄漏
* 与理论公式一致

**通过标准**

* 可视化检查通过
* 单元测试通过

**Git**

* 分支：`feat/m2-background-field`

---

### M3. 因子化 Eikonal 原型

**目标**
实现 `alpha` 的求解器原型。

**输出**

* FSM/FMM 原型
* `alpha`
* `tau = tau0 * alpha`

**检查项**

* 震源点初始化正确
* 扫描顺序稳定
* 更新单调
* 无明显假波纹

**通过标准**

* 常速介质中 `alpha ≈ 1`
* 平滑变速介质中 `tau` 连续
* 无 NaN

**Git**

* 分支：`feat/m3-factored-eikonal`

---

### M4. 安全导数重组

**目标**
构造 `grad_tau` 和 `lap_tau`，禁止危险直接差分。

**输出**

* `grad_alpha`
* `lap_alpha`
* `grad_tau`
* `lap_tau`

**检查项**

* 不直接对 `tau` 生硬做二阶差分
* 使用链式法则重组
* 震源附近数值可控

**通过标准**

* 可视化无星芒噪声
* 无异常尖峰扩散

**Git**

* 分支：`feat/m4-tau-derivatives`

---

### M5. PML 与差分引擎

**目标**
实现 `sigma`、`gamma`、`gamma'` 和固定卷积核微分。

**输出**

* `sigma_x/y`
* `gamma_x/y`
* `diff_ops.py`
* PML 拉普拉斯

**检查项**

* 不用 autograd 求二阶空间导数
* `gamma'` 用解析式
* PML 层只在边界起作用

**通过标准**

* 平面波进入边界后衰减
* 无明显边界反射

**Git**

* 分支：`feat/m5-pml-diffops`

---

### M6. 等效体源与 source mask

**目标**
实现 `RHS_eq` 与 punctured mask。

**输出**

* `RHS_eq`
* `loss_mask`
* source neighborhood 规则

**检查项**

* 震源点不出现 `0 * inf`
* PML 内 RHS 为 0
* mask 半径固定且写进配置

**通过标准**

* 数值稳定
* 可视化连续
* 无 NaN

**Git**

* 分支：`feat/m6-rhs-mask`

---

### M7. 网络骨架

**目标**
只实现前向网络，不训练。

**输出**

* `SpectralBlock`
* `NSNO2D`
* 输入输出接口

**检查项**

* 通道维度统一
* 实虚双通道一致
* 前向不爆显存

**通过标准**

* 前向可跑通
* 输出 shape 正确

**Git**

* 分支：`feat/m7-model-skeleton`

---

### M8. PDE Residual 前向检查

**目标**
把预处理、网络、残差拼起来，只做 forward residual。

**输出**

* `ResidualComputer`
* `loss_pde`
* debug plot

**检查项**

* 物理域内残差有限
* PML 行为合理
* 震源附近 mask 生效

**通过标准**

* 全图无 NaN/Inf
* 小样例 forward 成功

**Git**

* 分支：`feat/m8-residual-forward`

---

### M9. 最小训练实验

**目标**
先做最小 overfit，不做大训练。

**输出**

* 1 个小样本 overfit
* loss 曲线
* reconstruction 图

**检查项**

* loss 能下降
* 不出现数值爆炸
* 输出图像合理

**通过标准**

* 小样本可拟合
* 训练稳定

**Git**

* 分支：`feat/m9-min-train`

---

## 6. 每一步都要写“验收标准”

这是你文档里必须有的。
否则“看起来差不多”会毁掉整个节奏。

建议固定一个模板：

### 每个里程碑的验收模板

* 功能完成
* 单元测试通过
* 可视化检查通过
* 无 NaN/Inf
* 文档已更新
* 配置文件已更新
* 至少 1 个 debug case 已保存
* commit message 清晰
* 可以被他人在新环境复现

---

## 7. 写 Git 规则

这部分不要复杂，但一定要固定。

## Git Workflow

### 分支命名

* `main`：稳定主线
* `dev`：开发集成
* `feat/m1-grid-source`
* `feat/m2-background-field`
* `fix/pml-gamma`
* `exp/debug-source-mask`

### 提交规范

建议：

* `init: bootstrap repository`
* `feat: add 2d grid and point source`
* `feat: implement background field u0`
* `feat: add factored eikonal prototype`
* `fix: stabilize source mask near singularity`
* `refactor: split residual operator and diff ops`
* `docs: update DEV_PLAN for M4`

### 合并规则

* 不直接往 `main` 推
* 每个 milestone 单独分支
* milestone 通过验收后再 merge
* merge 前必须补文档
* 未通过测试不得合并

---

## 8. 写“失败回退规则”

这部分很多人不写，但非常有用。

## 回退规则

* 如果某里程碑连续两次修改后仍不稳定，停止继续叠加功能
* 回到上一个稳定 tag
* 在 `docs/CHANGELOG.md` 记录失败原因
* 若问题来自理论未冻结，先改 `DEV_PLAN.md`，再开新分支重做

建议每个里程碑过后打 tag：

* `v0.1-m0-bootstrap`
* `v0.1-m1-grid`
* `v0.1-m2-background`
* …

这样你永远能回去。

---

## 9. 再加一节“当前未决问题”

你现在这个项目很适合写这一节。
例如：

## Open Questions / 当前未决问题

1. `alpha` 求解器最终选 FSM 还是 FMM
2. source mask 半径取 `1.0h / 1.5h / 2.0h`
3. `lap_tau` 的离散实现最终采用哪一个安全版本
4. PML 厚度与 `sigma_max` 默认值
5. 复数张量是否全程双通道，还是部分模块用 complex dtype
6. NSNO 第一版是否直接上 FNO block

这能防止你把“尚未决定的问题”偷偷写进代码里。

---

# 最后给你一个最实用的建议

你这份文档最好不是纯说明文，而是**半规范、半清单**。
也就是说，每个 milestone 后面都留一个可勾选区：

```md
## M3. 因子化 Eikonal 原型
- [ ] Grid 接口稳定
- [ ] tau0 实现完成
- [ ] alpha 初始化完成
- [ ] FSM 一次扫描完成
- [ ] 四方向扫描完成
- [ ] 常速介质测试通过
- [ ] 变速介质测试通过
- [ ] 无 NaN/Inf
- [ ] 文档更新
- [ ] 合并到 dev
```

这比写长篇理论有用得多。

